% 《大学物理简明教程》第6章解答题6.3.1--6.3.9。
% 用户扫描印刷/PDF第156--157页；题意转述与独立解析，待图示和学术二校。
\subsection{解答题：第一定律、热容和实际热机}

\sourceproblem{6.3.1}{156}{156}
\problemtext 图6.3.1中系统从$a$沿曲线$acb$到$b$，从外界吸热350 J、对外做功126 J。
(1) 若改沿折线路径$adb$从$a$到$b$，气体对外做功42 J，求传入系统的热量。
(2) 从$b$沿曲线$ba$返回$a$，外界对系统做功84 J，判断吸放热并计算。
\workings
以系统对外做功$W>0$、系统吸热$Q>0$。
第一条路径由$\Delta U=Q-W$给
\[
U_b-U_a=350-126=\boxed{224\,\mathrm J}.
\]
(1) 状态函数内能差不随路径变，
\[
\boxed{Q_{adb}=\Delta U+W_{adb}=224+42=266\,\mathrm J}.
\]
(2) 反方向内能改变量$\Delta U_{ba}=-224$ J，外界对系统做功84 J意味着$W_{ba}=-84$ J，
\[
\boxed{Q_{ba}=\Delta U_{ba}+W_{ba}
=-224-84=-308\,\mathrm J}.
\]
故系统向外放热308 J。
\crosscheck 循环$acb+ba$的净功为$126-84=42$ J，净热也为$350-308=42$ J，满足完整循环$\Delta U=0$。

\sourceproblem{6.3.2}{156}{156}
\problemtext 用桨叶搅拌密闭容器内的液体，桨叶向系统输入功率2.24 kW，同时容器以0.586 kW的速率向环境散热。求系统每秒内能改变多少。
\workings
液体与容器为系统；桨叶输入的机械功率为正$P_{\mathrm{on}}=2.24$ kW，
系统放热$P_Q=-0.586$ kW。第一定律的功输入形式：
\[
\frac{dU}{dt}=\dot Q+\dot W_{\mathrm{on}}
=-0.586+2.24
=\boxed{1.654\,\mathrm{kW}}.
\]
若过程持续1 s，则$\boxed{\Delta U=1.654\times10^3\,\mathrm J}$。
\crosscheck 不可将桨叶输入机械功作为由系统\emph{对外}做功而用错符号；内能应增加。

\sourceproblem{6.3.3}{156--157}{156--157}
\problemtext 1 mol单原子理想气体由$300$ K升至$350$ K，分别经等容和等压过程，求吸热量、内能增加量与对外做功。
\workings
$\nu=1$ mol、$\Delta T=50$ K，单原子理想气体
$C_{V,m}=\frac32R$，$C_{P,m}=\frac52R$。
共通的内能改变：
\[
\boxed{\Delta U=\frac32R\Delta T
=\frac32(8.31446)(50)\approx623.6\,\mathrm J}.
\]
等容时$W_V=0$，$Q_V=\Delta U$，故
$\boxed{Q_V\approx623.6\,\mathrm J,\ W_V=0}$。
等压时$W_P=R\Delta T\approx415.7$ J，
\[
\boxed{Q_P=\frac52R\Delta T\approx1039.3\,\mathrm J}.
\]
\crosscheck 等压吸热多于等容过程，是因为额外的$R\Delta T$转为外界体积功；
两过程$\Delta U$相同。

\sourceproblem{6.3.4}{157}{157}
\problemtext 1 mol理想气体等压加热，温度升高$72$ K，吸热
$Q=1.60\times10^3$ J。求气体对外做功、内能增量和摩尔热容比$\gamma$。
\workings
等压功
\[
\boxed{W=\nu R\Delta T=R(72)\approx598.6\,\mathrm J}.
\]
根据第一定律
\[
\boxed{\Delta U=Q-W\approx1001.4\,\mathrm J}.
\]
由$C_{P,m}=Q/\Delta T=1600/72\,\mathrm{J/(mol\,K)}$，
$C_{V,m}=C_{P,m}-R$，故
\[
\boxed{\gamma=\frac{C_P}{C_V}
=\frac{1600/72}{1600/72-R}
\approx1.598}.
\]
\crosscheck $\gamma>1$；气体的题设等效自由度不必恰好是单原子或刚性双原子的整数模型。

\sourceproblem{6.3.5}{157}{157}
\problemtext 1 mol刚性多原子分子理想气体，初态$p_1=1$ atm、$T_1=27^\circ\mathrm C$；
准静态绝热压缩至$p_2=16$ atm。按刚性非线性多原子分子取$i=6$，求内能增量、对外做功及末态分子数密度。
\workings
刚性非线性多原子分子有3个平动和3个转动自由度，
$C_{V,m}=3R$、$\gamma=C_P/C_V=4/3$。
可逆绝热关系给
\[
T_2=T_1\left(\frac{p_2}{p_1}\right)^{(\gamma-1)/\gamma}
=(300.15\,\mathrm K)\,16^{1/4}
\approx600.3\,\mathrm K.
\]
所以
\[
\boxed{\Delta U=3R(T_2-T_1)\approx7.49\times10^3\,\mathrm J},
\]
$Q=0$意味着
$\boxed{W_{\mathrm{gas}}=-\Delta U\approx-7.49\times10^3\,\mathrm J}$。
末态分子数密度
\[
\boxed{n_2=\frac{p_2}{k_{\mathrm B}T_2}
\approx1.96\times10^{26}\,\mathrm{m^{-3}}}.
\]
\crosscheck 原书明确是“刚性多原子分子气体”而不是刚性双原子；$i=6,\gamma=4/3$。
若将多原子误读为双原子（$\gamma=7/5$），则末温度和内能结论会不同。

\sourceproblem{6.3.6}{157}{157}
\problemtext 向标准状态下2 mol氢气（刚性双原子）传热500 J，
分别求等容、等温、等压过程中终态的温度、压强和体积相对初态的变化。
\workings
取初始$T_0=273.15$ K、$p_0=1$ atm、$\nu=2$ mol，
$C_{V,m}=\frac52R$、$C_{P,m}=\frac72R$。
\emph{等容：}
\[
\Delta T_V=\frac{Q}{\nu C_V}=\frac{500}{5R}\approx12.03\,\mathrm K,
\]
\[
\boxed{T_V\approx285.18\,\mathrm K,\quad
\frac{p_V}{p_0}=\frac{T_V}{T_0}\approx1.0440,\quad
V_V=V_0}.
\]
\emph{等温：} $\Delta U=0$，$W=Q=500$ J，
\[
\nu RT_0\ln(V_T/V_0)=500
\quad\Longrightarrow\quad
\boxed{T_T=T_0,\quad
\frac{V_T}{V_0}=
\exp\left(\frac{500}{2RT_0}\right)\approx1.1164,
\quad\frac{p_T}{p_0}\approx0.8958}.
\]
\emph{等压：}
\[
\Delta T_P=\frac{500}{7R}\approx8.591\,\mathrm K,
\]
\[
\boxed{T_P\approx281.74\,\mathrm K,\quad
p_P=p_0,\quad
\frac{V_P}{V_0}=\frac{T_P}{T_0}\approx1.0315}.
\]
\crosscheck 所有过程吸热$Q>0$，等容无体积功、等温全部吸热转为功、
等压则同时升温并对外做功。

\sourceproblem{6.3.7}{157}{157}
\problemtext 一定量理想气体沿等压过程对外做功200 J。分别按单原子和刚性双原子气体模型，求过程所需的吸热量。
\workings
等压过程$W=\nu R\Delta T=200$ J，
单原子$C_P=\frac52R$，双原子$C_P=\frac72R$，
所以
\[
\boxed{Q_{\mathrm{mono}}=\frac52W=500\,\mathrm J},\qquad
\boxed{Q_{\mathrm{diatom}}=\frac72W=700\,\mathrm J}.
\]
\crosscheck 与同章填空题6.2.4的数值一致，但因原书分别给定题号，仍独立保留两份研究性解析。

\sourceproblem{6.3.8}{157}{157}
\problemtext 某理想气体的准静态过程方程为$Vp^{1/2}=a$，
$a$为正常量。若体积从$V_1$膨胀到$V_2$，求气体对外做的功。
\workings
由$Vp^{1/2}=a$得$p=a^2/V^2$，因此
\[
W=\int_{V_1}^{V_2}p\,dV
=a^2\int_{V_1}^{V_2}V^{-2}\,dV
=\boxed{a^2\left(\frac1{V_1}-\frac1{V_2}\right)}.
\]
也可按初态$p_1=a^2/V_1^2$改写为
$\boxed{W=p_1V_1(1-V_1/V_2)}$。
\crosscheck $V_2>V_1$时$W>0$；式中$a^2/V$的单位为能量。

\sourceproblem{6.3.9}{157}{157}
\problemtext 一台卡诺热机工作于$T_h=480$ K、$T_c=300$ K之间。
实际每从高温热源吸收1 kcal热量仅产生1.2 kJ机械能，比较实际效率与可逆卡诺效率上限。
\workings
取$1\,\mathrm{kcal}=4184\,\mathrm J$，则
\[
\boxed{\eta_{\mathrm{act}}=\frac{1200}{4184}
\approx0.2868}.
\]
同温热源的卡诺上限
\[
\boxed{\eta_C=1-\frac{300}{480}=0.375}.
\]
两者比值
\[
\boxed{\eta_{\mathrm{act}}/\eta_C\approx0.765}.
\]
\crosscheck 严格说题干的“卡诺热机”若指\emph{可逆卡诺循环}，
其效率必须等于$0.375$；现给出的实际输出仅约为上限的76.5\%，
因此该装置是\emph{在相同温度热源间工作的实际非可逆热机}，
原题措辞需要二校。
